% Sección (vi) Conclusión. Se inserta en main.tex después de Resultados y discusión:
%   \clearpage
%   % Sección (vi) Conclusión. Se inserta en main.tex después de Resultados y discusión:
%   \clearpage
%   % Sección (vi) Conclusión. Se inserta en main.tex después de Resultados y discusión:
%   \clearpage
%   % Sección (vi) Conclusión. Se inserta en main.tex después de Resultados y discusión:
%   \clearpage
%   \input{conclusion}
% No introduce cifras de resultados ni referencias nuevas: resume la sección anterior.
\section{Conclusión}
\label{sec:conclusion}\label{con-ini}

Este proyecto preguntó en qué condiciones D-MAB mejora Tabu Search frente a la
selección uniforme al programar cirugías en dos etapas. La literatura no había
medido ese aporte en un Tabu con espera estricta y reparación. Nuestra respuesta
es que no identificamos esa condición para D-MAB completo. Bajo congestión, el
resultado es no concluyente, porque el intervalo de la diferencia incluye el cero.

Sin embargo, el estudio deja tres hallazgos. Primero, la cartera supera al Tabu
prolongado con igual presupuesto. Segundo, esa ganancia es compatible con repetir
la búsqueda, porque no distinguimos la cartera de dos Tabu independientes.
Tercero, el reinicio Page--Hinkley no ayuda bajo congestión. De hecho, UCB1 sin
reinicio supera a D-MAB y, en una confirmación pre-registrada con semillas
nuevas, también a Tabu uniforme.

El aporte principal es metodológico. Diseñamos controles a igual presupuesto que
separan el aprendizaje del cómputo y de la repetición de la búsqueda. Además, dos
protocolos registrados antes de ejecutar restringieron nuestras decisiones. Así,
no declaramos un efecto que el intervalo no respalda, y confirmamos una señal
antes de afirmarla. Por último, generamos un catálogo de congestión que tensiona
la espera entre etapas.

Tres limitaciones acotan estas conclusiones. Primero, las semillas nuevas sobre
un mismo catálogo sintético no miden generalización a otros hospitales. Segundo,
la variabilidad observada dejó subdimensionada la comparación principal. Tercero,
comparamos por evaluaciones y en una sola máquina, no por tiempo.

Por tanto, el trabajo futuro debería contrastar cinco hipótesis. Primera: un
umbral $\lambda$ mayor reduce los reinicios y acerca D-MAB a UCB1. Segunda: un
crédito por rangos o por valores extremos entrega señal aunque las mejoras sean
escasas. Tercera: recompensar el movimiento aplicado, y no el propuesto, corrige
la distorsión de la reparación. Cuarta: reiniciar la búsqueda ante el
estancamiento reproduce la ganancia de la cartera. Quinta: un modelo sustituto
que prefiltre vecinos ahorra evaluaciones estrictas. Además, falta evaluar una
cartera con UCB1 frente a dos Tabu independientes. También faltan instancias
reservadas de 50 a 150 cirugías y una comparación a igual tiempo.

En conjunto, el proyecto muestra que el valor de un híbrido no depende de sumar
componentes de aprendizaje. Depende de comprobar, con controles, cuál de ellos
aporta. En este caso, aportó el componente más simple.
\label{con-fin}

% No introduce cifras de resultados ni referencias nuevas: resume la sección anterior.
\section{Conclusión}
\label{sec:conclusion}\label{con-ini}

Este proyecto preguntó en qué condiciones D-MAB mejora Tabu Search frente a la
selección uniforme al programar cirugías en dos etapas. La literatura no había
medido ese aporte en un Tabu con espera estricta y reparación. Nuestra respuesta
es que no identificamos esa condición para D-MAB completo. Bajo congestión, el
resultado es no concluyente, porque el intervalo de la diferencia incluye el cero.

Sin embargo, el estudio deja tres hallazgos. Primero, la cartera supera al Tabu
prolongado con igual presupuesto. Segundo, esa ganancia es compatible con repetir
la búsqueda, porque no distinguimos la cartera de dos Tabu independientes.
Tercero, el reinicio Page--Hinkley no ayuda bajo congestión. De hecho, UCB1 sin
reinicio supera a D-MAB y, en una confirmación pre-registrada con semillas
nuevas, también a Tabu uniforme.

El aporte principal es metodológico. Diseñamos controles a igual presupuesto que
separan el aprendizaje del cómputo y de la repetición de la búsqueda. Además, dos
protocolos registrados antes de ejecutar restringieron nuestras decisiones. Así,
no declaramos un efecto que el intervalo no respalda, y confirmamos una señal
antes de afirmarla. Por último, generamos un catálogo de congestión que tensiona
la espera entre etapas.

Tres limitaciones acotan estas conclusiones. Primero, las semillas nuevas sobre
un mismo catálogo sintético no miden generalización a otros hospitales. Segundo,
la variabilidad observada dejó subdimensionada la comparación principal. Tercero,
comparamos por evaluaciones y en una sola máquina, no por tiempo.

Por tanto, el trabajo futuro debería contrastar cinco hipótesis. Primera: un
umbral $\lambda$ mayor reduce los reinicios y acerca D-MAB a UCB1. Segunda: un
crédito por rangos o por valores extremos entrega señal aunque las mejoras sean
escasas. Tercera: recompensar el movimiento aplicado, y no el propuesto, corrige
la distorsión de la reparación. Cuarta: reiniciar la búsqueda ante el
estancamiento reproduce la ganancia de la cartera. Quinta: un modelo sustituto
que prefiltre vecinos ahorra evaluaciones estrictas. Además, falta evaluar una
cartera con UCB1 frente a dos Tabu independientes. También faltan instancias
reservadas de 50 a 150 cirugías y una comparación a igual tiempo.

En conjunto, el proyecto muestra que el valor de un híbrido no depende de sumar
componentes de aprendizaje. Depende de comprobar, con controles, cuál de ellos
aporta. En este caso, aportó el componente más simple.
\label{con-fin}

% No introduce cifras de resultados ni referencias nuevas: resume la sección anterior.
\section{Conclusión}
\label{sec:conclusion}\label{con-ini}

Este proyecto preguntó en qué condiciones D-MAB mejora Tabu Search frente a la
selección uniforme al programar cirugías en dos etapas. La literatura no había
medido ese aporte en un Tabu con espera estricta y reparación. Nuestra respuesta
es que no identificamos esa condición para D-MAB completo. Bajo congestión, el
resultado es no concluyente, porque el intervalo de la diferencia incluye el cero.

Sin embargo, el estudio deja tres hallazgos. Primero, la cartera supera al Tabu
prolongado con igual presupuesto. Segundo, esa ganancia es compatible con repetir
la búsqueda, porque no distinguimos la cartera de dos Tabu independientes.
Tercero, el reinicio Page--Hinkley no ayuda bajo congestión. De hecho, UCB1 sin
reinicio supera a D-MAB y, en una confirmación pre-registrada con semillas
nuevas, también a Tabu uniforme.

El aporte principal es metodológico. Diseñamos controles a igual presupuesto que
separan el aprendizaje del cómputo y de la repetición de la búsqueda. Además, dos
protocolos registrados antes de ejecutar restringieron nuestras decisiones. Así,
no declaramos un efecto que el intervalo no respalda, y confirmamos una señal
antes de afirmarla. Por último, generamos un catálogo de congestión que tensiona
la espera entre etapas.

Tres limitaciones acotan estas conclusiones. Primero, las semillas nuevas sobre
un mismo catálogo sintético no miden generalización a otros hospitales. Segundo,
la variabilidad observada dejó subdimensionada la comparación principal. Tercero,
comparamos por evaluaciones y en una sola máquina, no por tiempo.

Por tanto, el trabajo futuro debería contrastar cinco hipótesis. Primera: un
umbral $\lambda$ mayor reduce los reinicios y acerca D-MAB a UCB1. Segunda: un
crédito por rangos o por valores extremos entrega señal aunque las mejoras sean
escasas. Tercera: recompensar el movimiento aplicado, y no el propuesto, corrige
la distorsión de la reparación. Cuarta: reiniciar la búsqueda ante el
estancamiento reproduce la ganancia de la cartera. Quinta: un modelo sustituto
que prefiltre vecinos ahorra evaluaciones estrictas. Además, falta evaluar una
cartera con UCB1 frente a dos Tabu independientes. También faltan instancias
reservadas de 50 a 150 cirugías y una comparación a igual tiempo.

En conjunto, el proyecto muestra que el valor de un híbrido no depende de sumar
componentes de aprendizaje. Depende de comprobar, con controles, cuál de ellos
aporta. En este caso, aportó el componente más simple.
\label{con-fin}

% No introduce cifras de resultados ni referencias nuevas: resume la sección anterior.
\section{Conclusión}
\label{sec:conclusion}\label{con-ini}

Este proyecto preguntó en qué condiciones D-MAB mejora Tabu Search frente a la
selección uniforme al programar cirugías en dos etapas. La literatura no había
medido ese aporte en un Tabu con espera estricta y reparación. Nuestra respuesta
es que no identificamos esa condición para D-MAB completo. Bajo congestión, el
resultado es no concluyente, porque el intervalo de la diferencia incluye el cero.

Sin embargo, el estudio deja tres hallazgos. Primero, la cartera supera al Tabu
prolongado con igual presupuesto. Segundo, esa ganancia es compatible con repetir
la búsqueda, porque no distinguimos la cartera de dos Tabu independientes.
Tercero, el reinicio Page--Hinkley no ayuda bajo congestión. De hecho, UCB1 sin
reinicio supera a D-MAB y, en una confirmación pre-registrada con semillas
nuevas, también a Tabu uniforme.

El aporte principal es metodológico. Diseñamos controles a igual presupuesto que
separan el aprendizaje del cómputo y de la repetición de la búsqueda. Además, dos
protocolos registrados antes de ejecutar restringieron nuestras decisiones. Así,
no declaramos un efecto que el intervalo no respalda, y confirmamos una señal
antes de afirmarla. Por último, generamos un catálogo de congestión que tensiona
la espera entre etapas.

Tres limitaciones acotan estas conclusiones. Primero, las semillas nuevas sobre
un mismo catálogo sintético no miden generalización a otros hospitales. Segundo,
la variabilidad observada dejó subdimensionada la comparación principal. Tercero,
comparamos por evaluaciones y en una sola máquina, no por tiempo.

Por tanto, el trabajo futuro debería contrastar cinco hipótesis. Primera: un
umbral $\lambda$ mayor reduce los reinicios y acerca D-MAB a UCB1. Segunda: un
crédito por rangos o por valores extremos entrega señal aunque las mejoras sean
escasas. Tercera: recompensar el movimiento aplicado, y no el propuesto, corrige
la distorsión de la reparación. Cuarta: reiniciar la búsqueda ante el
estancamiento reproduce la ganancia de la cartera. Quinta: un modelo sustituto
que prefiltre vecinos ahorra evaluaciones estrictas. Además, falta evaluar una
cartera con UCB1 frente a dos Tabu independientes. También faltan instancias
reservadas de 50 a 150 cirugías y una comparación a igual tiempo.

En conjunto, el proyecto muestra que el valor de un híbrido no depende de sumar
componentes de aprendizaje. Depende de comprobar, con controles, cuál de ellos
aporta. En este caso, aportó el componente más simple.
\label{con-fin}
